\section{Mecanismos de recompensa: reward intrínseca y Multi-Agent}

Esta sección aborda la reward. Yao Shunyu (姚顺雨) mencionó que una de las direcciones clave del desarrollo de los Agent es que tengan su propia reward y puedan explorar por sí mismos; otra dirección es Multi-Agent, en la que los agentes forman una estructura organizativa entre sí. Aquí la reward no es solo una fórmula de RL: es lo central para que el Agent obtenga retroalimentación efectiva del entorno, vaya más allá de imitar datos humanos y se mejore a sí mismo durante la tarea.

\begin{figure}[H]
\centering
\includegraphics[width=0.94\textwidth]{reward-multi-agent.png}
\caption{Dos direcciones clave para los Agent: tener su propia reward y que múltiples agentes formen una estructura organizativa. Diagrama conceptual propio, elaborado a partir de la conversación de 00:27:19--01:03:00.}
\end{figure}

\begin{knowledgebox}{Lectura de la figura: la reward y el multi-agent son dos vías de expansión distintas}
La reward permite que un solo Agent explore, pruebe, se equivoque y optimice; Multi-Agent permite que varios Agent formen organizaciones más complejas mediante la división del trabajo, la colaboración y la competencia. La primera resuelve la señal de aprendizaje; la segunda, la estructura organizativa.
\end{knowledgebox}

\subsection{Dónde está la dificultad de la reward}

La sección anterior separó la reward y el multi-agent en dos vías de expansión; esta sección examina primero por qué la reward es difícil. En las tareas del mundo real, la reward es difícil de definir. Los juegos de tablero tienen un ganador y un perdedor claros, las tareas de código tienen pruebas y las tareas web pueden tener un estado de finalización; pero los objetivos de muchas tareas reales son de largo plazo, dispersos, subjetivos o no completamente observables. Si la reward está mal diseñada, el Agent puede optimizar el objetivo equivocado o aprovechar los huecos de la evaluación.

\begin{figure}[H]
\centering
\includegraphics[width=0.96\textwidth]{reward-definition-problem.png}
\caption{El problema de definir la recompensa: en las tareas reales, la reward es difícil porque es de largo plazo, dispersa y no completamente observable. Diagrama conceptual propio, elaborado a partir de la conversación de 00:40:00--01:03:00.}
\end{figure}

\begin{warningbox}{El riesgo del reward hacking}
Si la reward solo premia el resultado superficial, el Agent puede aprender a hacer trampa. Por ejemplo, si solo se premia la tasa de clics, puede fabricar titulares sensacionalistas; si solo se premia completar la tarea, puede saltarse los permisos; si solo se premia pasar las pruebas, puede escribir código frágil. Las tareas empresariales y del mundo real requieren en particular auditoría y restricciones de seguridad.
\end{warningbox}

\begin{importantbox}{La fórmula intuitiva de la reward}
\[
\text{learning signal} = r(s_t, a_t, s_{t+1})
\]
Donde \(s_t\) es el estado antes de actuar, \(a_t\) es la acción del Agent, \(s_{t+1}\) es el estado después de actuar y \(r\) es la función de recompensa que evalúa ese cambio. La dificultad de las tareas reales está en que el estado es incompleto, las consecuencias de las acciones llegan con retraso y los criterios de evaluación suelen definirlos en conjunto las personas y las organizaciones.
\end{importantbox}

\subsection{Multi-Agent: de la capacidad individual a la estructura organizativa}

Lo central de Multi-Agent no es «abrir varias instancias del modelo», sino que los Agent formen entre sí división del trabajo, comunicación, colaboración, competencia y una estructura organizativa. Esto se parece a una organización humana y también trae problemas nuevos: quién asigna las tareas, quién verifica los resultados, quién asume los conflictos y cómo evitar las alucinaciones colectivas.

\begin{knowledgebox}{Términos clave: Reward y Multi-Agent}
\begin{tabularx}{\textwidth}{p{0.22\textwidth}p{0.32\textwidth}X}
\toprule
Concepto & Problema que resuelve & Riesgos comunes \\
\midrule
External reward & Evaluación externa que da el entorno o una persona & Costosa, dispersa y propensa al sobreajuste. \\
Intrinsic reward & El Agent genera su propia señal de exploración & Puede desviarse de los objetivos humanos. \\
Multi-Agent & Varios Agent dividen el trabajo y colaboran & Costo de comunicación, conflictos y límites de responsabilidad. \\
Verification & Verificar si la salida del Agent es confiable & El propio verificador también puede equivocarse. \\
\bottomrule
\end{tabularx}
\end{knowledgebox}

La dificultad de la reward y el multi-agent también está en que cambian el «carácter social» del sistema. Cuando un solo Agent se equivoca, el problema puede rastrearse hasta una sola política; cuando varios Agent colaboran, el error puede venir de la asignación de tareas, de malentendidos en la comunicación, de óptimos locales, de defectos del verificador o de incentivos desalineados. Cuanto más se acerca a una organización real, más necesita el sistema de Agent una capa de gobernanza, y no solo un modelo más potente.

\begin{warningbox}{Más agentes no significa automáticamente más capacidad}
Un sistema de múltiples Agent puede aportar división del trabajo y paralelismo, pero también puede traer mayores costos de comunicación, modos de falla más complejos y responsabilidades poco claras. Solo cuando la tarea es descomponible por naturaleza, el mecanismo de verificación es claro y el protocolo de colaboración es estable, es más probable que un sistema multiagente supere a un solo Agent.
\end{warningbox}

\subsection{Resumen de la sección}

La reward determina qué aprende el Agent; Multi-Agent determina cómo se organizan los Agent. La dificultad de la segunda mitad no es que el modelo converse mejor, sino que el sistema pueda explorar y colaborar con seguridad en entornos reales.
